\section{视频生成的进步，不只是“画面更清楚”}

这场课从几段生成视频开始，但真正的问题不是挑选哪一帧最漂亮，而是理解 2022--2024 年间为什么运动、反射、镜头语言、文本遵循和可编辑性同时进步。讲者给出的主线很克制：把媒体压缩为可处理的 token 序列，用 Flow Matching 训练一个大而统一的 Transformer，再让数据、计算与模型规模共同增长。后文将把这句话拆成可核算的表示、学习目标、架构、数据和评估系统。

\subsection{三种开场样例分别在证明什么}

开场样例承担三种不同证据角色：冲浪考拉展示复杂主体与运动，幽灵镜像展示几何与反射一致性，跑步者编辑展示输入视频的内容保持和指令控制。本节先建立这套读图框架：先问“任务是什么”，再问“哪些属性保持了”，最后问“这是不是稳定能力”；否则几个成功片段很容易被误写成完整可靠性，也无法为后面的架构判断提供证据。

\lecturefigure{slide-064-moviegen-demo-koala.jpg}{文本到视频样例：主体、相机、海浪与冲浪动作需要共同保持}{00:01:50}

冲浪片段的难点不是生成一只考拉，而是让主体外观、冲浪板、海浪、镜头距离和连续运动在多个帧中一致。单帧图像模型只需优化空间 plausibility；视频模型还必须控制时间轴上的身份稳定、运动连续和遮挡恢复。

\lecturefigure{slide-091-moviegen-demo-ghost.jpg}{镜像样例：反射、视角和主体身份形成跨区域约束}{00:02:22}

镜像场景要求模型把“真实主体”和“镜中主体”绑定成同一对象，同时遵守视角变化。它支持模型已经学到某些视觉规律，却不证明模型拥有显式光学引擎；更可靠的测试需要改变镜面角度、遮挡和主体位置，观察规律是否系统泛化。

\lecturefigure{slide-112-moviegen-editing-intro.jpg}{指令式视频编辑：同一输入分别增加物体、替换服装与改变环境}{00:02:48}

编辑任务的评价维度与纯生成不同。除了输出是否好看，还要检查未编辑区域是否保持、目标区域是否严格执行指令、动作节奏是否延续，以及修改是否在整段视频中持续存在。后文会看到，这类能力并非直接从基础 T2V 模型自动出现，而需要额外训练路径。

\begin{warningbox}{成功样例不是分布级可靠性}
视频演示能证明“至少存在一个成功输出”，不能证明对 prompt 分布、随机种子、长尾对象和复杂交互都稳定。可信评估需要固定 prompt 集、公开失败样本、多个随机种子和人类成对比较；若只展示最佳样例，就无法估计成功率。
\end{warningbox}

\teachervoice{00:01:13--00:03:42，老师用三个开场 demo 区分生成质量、复杂反射和指令编辑。他强调近两年的提升不是单帧美学，而是模型开始处理运动、关系与控制。}

\subsection{从 Make-A-Video 到 Movie Gen：什么真正改变了}

下面两张图把时间线和讲座结论放在一起。第一张比较 2022 与 2024 的可见差异；第二张把改进归结为统一架构与规模化，而不是某个孤立 trick。理解这段历史时，应把算法、数据、硬件、训练基础设施和评估方法视为共同变量。

\lecturefigure{slide-150-video-generation-progress.jpg}{两年进步的直观对照：Make-A-Video 与 Movie Gen}{00:03:35}

左侧早期模型已经能生成动态内容，但主体边界、运动稳定和细节连续性明显更弱；右侧 Movie Gen 在复杂海浪、相机跟随和主体一致性上更强。图支持“能力前沿显著移动”，但无法从视觉对照单独分解增益来自架构还是训练规模。

\lecturefigure{slide-155-moviegen-paper-thesis.jpg}{讲座核心命题：规模化数据、计算与参数的简单 Transformer 也适用于视频}{00:05:02}

这里的“简单”是相对 U-Net 级联和高度专用视频模块而言，并不意味着系统便宜。Movie Gen 仍需要 TAE、多个文本编码器、Flow Matching、分布式并行、数据清洗、渐进训练和后训练。核心简化发生在 backbone：尽量复用可预测扩展的 Transformer，而不是为每个分辨率和模态重新发明网络。

\begin{importantbox}{架构统一的工程含义}
统一架构不是“所有模态完全相同”，而是把模态专属复杂度集中在 encoder、decoder 与 conditioning interface，让主干继续使用可规模化的 Transformer blocks。这样模型、kernel、并行策略和训练基础设施能够复用，扩展方向也更清晰。
\end{importantbox}

\subsection{历史路线：专用小模型走向统一大模型}

前面的样例展示结果，本节转向机制史。时间线要回答两个问题：早期方法依靠哪些专用结构，为什么 2024 年的大模型更偏好统一 Transformer。读图时不要把“新”自动等同于“优”；关键是训练规模和基础设施改变后，哪种偏置更容易持续获益。

\lecturefigure{slide-174-video-generation-history.jpg}{视频生成路线：扩散、U-Net 级联与 Transformer/Flow Matching 的统一化}{00:09:39}

早期 Video Diffusion、Make-A-Video 和 Imagen Video 常把 U-Net、扩散、插帧与多级超分组合成复杂 pipeline。到 Sora、Movie Gen 等路线，研究重心转向统一 token backbone、Flow Matching 和更大规模训练。变化不是所有旧模块消失，而是系统开始优先减少架构分叉。

\lecturefigure{slide-175-moviegen-overview.jpg}{Movie Gen Video 的规模、任务与训练数据概览}{00:10:37}

Movie Gen Video 是约 30B 参数的 joint text-to-image/text-to-video foundation model，最大上下文约 73K video tokens，对应最长 16 秒、16 fps 的视频训练设置。讲者用 $O(100\mathrm{M})$ 视频和 $O(1\mathrm{B})$ 图像描述数据量级；这些是规模级口径，不应被解读为精确公开计数。

\begin{knowledgebox}{为什么联合训练图像与视频}
图像数据数量更多、语义覆盖更广，可帮助模型学习对象、风格和构图；视频数据提供运动与时间关系。联合训练让同一 backbone 同时吸收空间和时空规律，也允许图像被视为单帧视频进入统一表示。
\end{knowledgebox}

\teachervoice{00:07:57--00:11:19，老师把路线概括为从“小规模、专用架构”转向“架构统一、简单 Transformer 和规模化”。他同时给出 30B、73K tokens、16 秒和 16 fps 的实际 workload，提醒统一不等于廉价。}

\subsection{本章小结}

视频生成的进步包含空间质量、时间一致性、复杂关系和可控编辑。Movie Gen 的核心命题是用可扩展 Transformer 统一 backbone，但完整系统仍依赖表示压缩、学习目标、数据、并行与评估。样例是能力证据，不是可靠性或因果世界模型的充分证明。

